\begin{table*}[t]\centering\small
\caption{Relationship to existing evaluation guidance. These are field correspondences, not an empirical comparison of checklists. Arp P-numbers refer to Section~2 of the published paper~\cite{arp2022dos}.}\label{tab:guideline-crosswalk}
\setlength{\tabcolsep}{4pt}
\begin{tabularx}{\textwidth}{@{}P{.14\textwidth}P{.37\textwidth}Y@{}}\toprule
Entries & Established obligation & Agent-specific specification used here \\\midrule
C1, C4 & Realistic observer information and splits (Juarez et al.; Mathews et al.)~\cite{juarez2014critical,mathews2023critical}. & Separate transit, recipient and active access; identify supplied stage boundaries and association across model/tool links. \\\addlinespace[3pt]
C2, C3 & Sampling and label validity (Arp P1--P2). & Distinguish requested intent, logged action and checked outcome; state exclusions, null responses and recoverable errors. \\\addlinespace[3pt]
C4, C5 & Suitable comparisons (Arp P6), changed conditions and adaptive evaluation~\cite{juarez2014critical,mathews2023critical,tramer2020adaptive}. & Name composition, topology, scaffold or model changes; distinguish fixed algorithm from refitting on protected traffic. \\\addlinespace[3pt]
C6 & Protection overhead and realistic evaluation~\cite{dyer2012peekaboo,mathews2023critical}; deployment scope (Arp P9). & Bind original-task utility to executed cover, and bytes/latency to the same protection condition and population. \\\addlinespace[3pt]
C7 & Application-relevant metrics and prevalence (Arp P7--P8). & Report event detection separately from supplied-unit classification; attach uncertainty and base rates to the named monitoring target. \\\addlinespace[3pt]
C8 & Evidence record supporting the correspondences above. & Give source version, result locator and accessible artifacts; distinguish source inspection from experimental reproduction. \\
\bottomrule\end{tabularx}\end{table*}
